% !PRINTED 660
\weberpage{660}

The theorem will be proved if it can be shown of any power
$\varphi^k$ of $\varphi$ whose exponent is not divisible by $q$
that it belongs to the principal class. For if $\alpha, \beta$
are numbers in $\Omega_m$ and
\[
\varphi^m = \varepsilon' \alpha, \quad \varphi^k = \varepsilon'' \beta,
\]
then we can determine the integral rational numbers $x, y$ in such
a way that
\[
mx + ky = 1
\]
holds, and then
\[
\varphi = \varphi^{mx} \varphi^{ky} = \varepsilon''' \alpha^x \beta^y,
\]
hence, if $\varepsilon', \varepsilon'', \varepsilon'''$ are units,
$\varphi$ itself is also associated with a number.

In the present state of the matter, the proof for an even $m$ must
be conducted along an entirely different path than for an odd one,
and so we occupy ourselves first with the easier case of odd
$m$ $^1$).

The assumptions made in a) about $\varphi_n$ can be expressed, if
$\alpha, \beta$ again denote any numbers of the field $\Omega_m$
and $\varepsilon$ unit functionals, by the two formulae:
\weq{1}{\varphi_n^m = \varepsilon \beta,}
\weq{2}{\varphi_n = \varepsilon \alpha \varphi_1^n.}
We shall now solve the problem step by step.

If first, as before, $\sigma$ is the prime factor contained in
$q$, which is a number existing in $\Omega_m$, and $\sigma^k$ the
power of $\sigma$ contained in $\varphi$, then we set
\weq{3}{\varphi = \sigma^k \psi,}
in which $\psi$ denotes a functional coprime to $q$ that satisfies
the same conditions (1), (2) as $\varphi$; and if one of the two
belongs to the principal class, the same holds of the other. The
theorem a) therefore need only be proved under the assumption that
$\varphi$ is coprime to $q$.

Now let $p$ be a prime different from $q$ which, decomposed into
prime factors, has the expression:
\weq{4}{p = \prod^{\xi} \mathfrak{p}_\xi,}
%CONT>>

\noindent
$^1$) Cf., moreover, the new proof by \gesperrt{Hilbert}, in which
the difference between even and odd $m$ recedes far more
(S. 648).
